
\section{Limitations, Future Work, and Broader Impact}
% \todo{rooflines and MFU validates, need to run this on a real simulator and finetune}
Our current formulation treats the neural network architecture as fixed.
A natural extension would co-optimize network architecture, hardware datapath, and compiler decisions in a single framework, building on prior work in hardware-aware NAS.
The LP relaxation may produce fractional solutions that require rounding; investigating tighter formulations or branch-and-bound within the LP would improve solution quality.
In our audits this looseness is quantified rather than hypothetical: rounding is essentially lossless on 9 of 12 workloads (relaxation-to-replay gaps of 0.00--0.64\%), while on the weight-heavy Llama-1 65B and Llama-3 70B the deployed schedule is certified optimal for its integral restriction by an exact chain dynamic program yet sits 23--30\% above the LP lower bound: certified relaxation looseness concentrated in the weight-residency rows, and the sharpest motivation for the tighter disaggregated epigraph of Section~\ref{sec:v1}.
A systematic account of when each rung of the V0$\to$V1$\to$V2$\to$V3 hierarchy pays for its added variables (V0 already models placement and rematerialization; V1's joint tiling--fusion coupling is what unlocks the free-$y$ gains) is deferred to the expanded analysis accompanying the V3 experiments.
The LLM-guided outer loop is preliminary; a systematic study of prompt engineering, LLM-Vizier integration strategies, and the interaction between LLM reasoning and LP dual information would strengthen this component.
Finally, extending the formulation to optimize for training workloads (which require backward-pass memory management and gradient accumulation) is an important direction.
As a concrete target for the decode-regime extension: DeepSeek-V3 inference is bottlenecked in the autoregressive decode phase by DRAM bandwidth spent reloading the mixture-of-experts weights (${\sim}98\%$ of decode traffic) because at serving batch the top-8 routing collides across tokens and forces nearly the entire 671B-parameter model to stream from memory every step; the KV cache (MLA-compressed) and compute are negligible by comparison.
The single largest real-world decision for this bottleneck, expert parallelism across devices, is inherently multi-node, whereas the co-design program here optimizes a single custom chip; on one chip the parallelization we can co-design is therefore the compute-side one (raising the speculative/parallel-decode width $K$ and provisioning enough compute to absorb it, while holding the hottest experts in on-chip memory), and the device-spreading answer belongs either to the multi-node compiler formulation or to a future heterogeneous, multi-chiplet extension that is itself a co-design axis.
More efficient accelerator design tools can reduce the energy consumption and carbon footprint of datacenter ML inference, which accounts for a growing share of global compute.
By lowering the engineering effort and deployment scale required for specialized accelerators to be economically viable, our work may also broaden access to high-performance ML inference beyond the largest technology companies.
